\bookchapter{Capítulo 63: Vistas, estado y navegación}

\chapterintro{Ya sabes hacer apps de una sola pantalla. Este capítulo resuelve dos problemas de «crecimiento»: cómo saltar de unas páginas a otras (navegación) y cómo fluyen los datos entre páginas (gestión de estado). Al terminarlo podrás montar el esqueleto clásico de una app multipágina: «lista detalle edición».}

\booksection{Panorama del flujo de datos para gestionar el estado}

@State no es la única herramienta de estado de SwiftUI. Mira primero este panorama para entender «cuál usar en cada momento»:

\begin{center}
\begin{tabularx}{\linewidth}{|X|X|X|X|}
\hline
\textbf{Herramienta} & \textbf{Qué gestiona} & \textbf{Quién la usa} & \textbf{Escenario típico} \tabularnewline
\hline
@State & Estado temporal privado de la página & La vista actual & Contadores, texto de un campo \tabularnewline
@Binding & Tomar prestado el estado de otro & Una vista hija modifica a la padre & Componentes personalizados \tabularnewline
@StateObject / @Observable & Lógica compartida entre páginas & Modelo de vista (ViewModel) & Carga de datos, lógica de negocio \tabularnewline
@EnvironmentObject / .environment & Estado global compartido & Toda la app & Sesión del usuario, ajustes del tema \tabularnewline
\hline
\end{tabularx}
\end{center}

\begin{center}\small\color{gray}Panorama de las herramientas de estado de SwiftUI (en 2026, @Observable ya es lo habitual)\end{center}

\begin{lstlisting}[language=Swift]
// la forma moderna de iOS 17+: la macro @Observable (recomendada para proyectos nuevos)
import Observation

@Observable
class CounterModel {                 // el view model
    var count = 0
    func increment() { count += 1 }
}

struct CounterView: View {
    @State private var model = CounterModel()   // mantiene el modelo

    var body: some View {
        VStack {
            Text("Has pulsado \(model.count) veces")
            Button("Púlsame") { model.increment() }
        }
    }
}
\end{lstlisting}

Consejo directo para principiantes: usa @State para lo privado de la página y clases @Observable para lo que cruza páginas o lleva lógica. El @StateObject/@ObservableObject de los tutoriales viejos es la generación anterior: basta con saber leerlo; el código nuevo, hazlo como aquí.

\booksection{NavigationStack: la «pila» para saltar de página}

Desde iOS 16, la navegación multipágina usa NavigationStack: funciona como una pila de platos, apilas uno sobre otro y al volver se van sacando:

\begin{lstlisting}[language=Swift]
struct HomeView: View {
    var body: some View {
        NavigationStack {
            List {
                NavigationLink("Las cuentas de Luis") {
                    DetailView(name: "Las cuentas de Luis")   // página destino
                }
                NavigationLink("Las cuentas de María") {
                    DetailView(name: "Las cuentas de María")
                }
            }
            .navigationTitle("Lista de cuentas")   // título grande arriba
        }
    }
}

struct DetailView: View {
    let name: String
    var body: some View {
        Text("Detalles de \(name)")
            .navigationTitle(name)
            .navigationBarTitleDisplayMode(.inline)
    }
}
\end{lstlisting}

NavigationLink es el «enlace que se toca para navegar»; navigationTitle genera automáticamente el título de la barra de navegación y el botón de retroceso. Para navegar por código (por ejemplo, volver solo tras pulsar «Guardar»), usa @Environment(\textbackslash{}.dismiss) para cerrar la página actual. Con estas dos técnicas cubres el 90 \% de las necesidades de navegación.

\booksection{Ventanas modales y pestañas: sheet y TabView}

Además de «empujar páginas», hay dos interacciones clásicas:

\begin{lstlisting}[language=Swift]
// 1. hoja que sube desde abajo, media pantalla o completa (la favorita de los formularios de alta)
struct HomeView: View {
    @State private var showAdd = false
    var body: some View {
        Button("Apuntar gasto") { showAdd = true }
            .sheet(isPresented: $showAdd) {
                AddRecordView()     // la pantalla de alta que se abre
            }
    }
}

// 2. barra de pestañas inferior (el "esqueleto" de la app)
struct MainTabView: View {
    var body: some View {
        TabView {
            RecordListView()
                .tabItem { Label("Cuentas", systemImage: "book") }
            StatsView()
                .tabItem { Label("Estadísticas", systemImage: "chart.pie") }
            SettingsView()
                .tabItem { Label("Ajustes", systemImage: "gear") }
        }
    }
}
\end{lstlisting}

\begin{tcbinfo}{Consejo: ¿cuándo usar cada una?}
Regla práctica: para el flujo principal, empuja páginas con NavigationStack (lista detalle); para acciones temporales, ventanas con sheet (añadir, editar, confirmar); para bloques grandes al mismo nivel, TabView (pestañas inferiores). La combinación de los tres es el esqueleto de casi todas las apps.

\end{tcbinfo}

\booksection{List: la lista por excelencia de iOS}

El componente más usado y con más «sabor Apple» de iOS es List: trae de fábrica separadores, borrado al deslizar y resaltado al tocar:

\begin{lstlisting}[language=Swift]
struct RecordListView: View {
    @State private var records = ["Comida 25.5", "Metro 6.0", "Nómina +8000"]

    var body: some View {
        List {
            ForEach(records, id: \.self) { r in
                Text(r)
            }
            .onDelete { indexSet in          // desliza a la izquierda para borrar
                records.remove(atOffsets: indexSet)
            }
        }
    }
}
\end{lstlisting}

Lista + NavigationLink + sheet + TabView: combinando estas cuatro piezas ya puedes montar el esqueleto completo de una app de «lista de cuentas abrir el detalle pulsar + y registrar». Así se construye la libreta de cuentas del próximo capítulo.

\begin{tcbpractice}{Manos a la obra: monta un «esqueleto de tres páginas»}
\booksection{Ciclo de vida de la app: primer plano y segundo plano}

Como el ciclo de vida de las Activity de Android, las apps de iOS también pasan por estados de «nacer, primer plano, segundo plano y morir». En SwiftUI se observan con el valor de entorno scenePhase:

\begin{lstlisting}[language=Swift]
@main
struct MyApp: App {
    @Environment(\.scenePhase) private var phase

    var body: some Scene {
        WindowGroup {
            ContentView()
        }
        .onChange(of: phase) { _, newPhase in
            switch newPhase {
            case .active:    print("vuelta a primer plano: reanudar el refresco")
            case .inactive:  print("en transición (llamada o el instante en que pasa a segundo plano)")
            case .background: print("pasa a segundo plano: guardar el borrador")
            @unknown default: break
            }
        }
    }
}
\end{lstlisting}

Usos típicos: refrescar los datos al volver al primer plano y guardar el borrador al pasar a segundo plano. iOS es muy estricto con el trabajo en segundo plano: una app normal se suspende (se congela) poco después de ir a segundo plano, así que no cuentes con colar trabajo pesado ahí. Si de verdad necesitas tareas programadas, usa las notificaciones locales del capítulo 68, no un «servicio residente».

Haz un esqueleto de libreta de cuentas con datos falsos: un TabView con tres pestañas (cuentas / estadísticas / ajustes) la página de cuentas muestra los datos falsos con una List y al tocar una fila salta al detalle (NavigationStack) el botón «+» de la esquina superior derecha abre un sheet de registro. Sin guardar datos y sin internet, pura práctica de esqueleto: si este esqueleto te sale bien, rellenarlo con datos reales en el capítulo 65 será coser y cantar.

\end{tcbpractice}

\begin{tcbsummary}{Resumen del capítulo}
\begin{itemize}
  \item Las cuatro herramientas de estado: @State (privado de la página), @Binding (prestado), @Observable (lógica) y .environment (global)
  \item En proyectos nuevos, escribe los view models con la macro @Observable; con el @ObservableObject de los tutoriales viejos basta con que sepas leerlo
  \item NavigationStack empuja páginas; @Environment(\textbackslash{}.dismiss) permite volver por código
  \item En el flujo principal se empujan páginas, las acciones temporales se abren con sheet y los bloques paralelos van en un TabView
  \item List es la lista por excelencia: separadores y borrado al deslizar, de fábrica
\end{itemize}

\end{tcbsummary}
